\textbf{结论名称：} 函数值域——单调性法

\textbf{适用条件：} 函数在指定区间上有定义且能判定其单调性（递增或递减），通常用于闭区间，也可借助极限推广至开区间。

\textbf{变量说明：} 设函数 \(y=f(x)\) 在区间 \(I=[a,b]\) 上单调，其中 \(a<b\) 为给定区间端点。

\textbf{核心公式：}
\[
\boxed{\text{值域}=
\begin{cases}
[f(a),\,f(b)], & f(x)\text{ 在 }[a,b]\text{ 上单调递增}\\[4pt]
[f(b),\,f(a)], & f(x)\text{ 在 }[a,b]\text{ 上单调递减}
\end{cases}}
\]

\textbf{使用场景：} 已用定义、图像或导数确定函数单调性后，快速写出闭区间上的值域，是求值域最基础、最高频的方法，常作为复合函数值域的中间步骤。

\textbf{简单例子：} 
求 \(f(x)=2x+1\) 在 \([1,3]\) 上的值域。
解：一次函数斜率 \(2>0\)，在 \(\mathbb{R}\) 上单调递增，故在 \([1,3]\) 上值域为 \([f(1),\,f(3)]=[3,\,7]\)。